%% kapa: exam=1 round=23 year=2012 period=2 form=A subject=appraisal_law subject_ko=부동산관계법규 seq=1 number=41
%% source: Q-Net 감정평가사 기출문제 (cst00309) artlSeq=5052580 file=raw/1008/5052580/2012년 제23회 감정평가사 제1차시험 기출문제(2교시 A형).pdf
%% conversion: kordoc 4.24.51 -> tools/gen1.py (auto); review=proofread-sonnet; flags=none
\begin{question}{41}
국토의 계획 및 이용에 관한 법령상 광역도시계획 및 도시기본계획에 관한 설명으로 옳지 않은 것은?
\begin{choices}
  \item 국토해양부장관은 광역도시계획을 수립하려면 미리 공청회를 열어 주민과 관계 전문가 등으로부터 의견을 들어야 하며, 공청회에서 제시된 의견이 타당하다고 인정하면 광역도시계획에 반영하여야 한다.
  \item 광역계획권이 같은 도의 관할 구역에 속하여 있는 경우 시장 또는 군수가 협의를 거쳐 요청하면 도지사가 단독으로 광역도시계획을 수립할 수 있다.
  \item 도시기본계획의 수립권자는 특별시장ㆍ광역시장ㆍ시장 또는 군수이며, 도시기본계획의 수립기준은 승인권자인 시ㆍ도지사가 정한다.
  \item 도지사는 시장 또는 군수가 수립한 도시기본계획을 승인하려면 관계 행정기관의 장과 협의한 후 지방도시계획위원회의 심의를 거쳐야 한다.
  \item 특별시장ㆍ광역시장ㆍ시장 또는 군수는 5년마다 관할 구역의 도시기본계획에 대하여 그 타당성 여부를 전반적으로 재검토하여 정비하여야 한다.
\end{choices}
\end{question}
\answer{3}
